\section{La media interior a corte positivo}\label{app:interior}
Para hallar la contribución interior a \eqref{eq:signed-mean} fijamos el corte \(b>0\).
Cuando quedan \(m>\lfloor bn\rfloor=j\) filas, escribimos \(x=m/n\).
La elección de una casilla modifica el predictor en orden \(1/m\) después de descontar
\(\log p_c\). Para identificar lo que queda en orden \(1/m^2\) empezamos en un padre regular,
antes de promediar sobre las historias.

\subsection{La inserción de una casilla}
Sean \(c\) la casilla elegida, \(\mathcal J=\familias(c)\) sus cuatro líneas, \(I\) las líneas retenidas,
\(F\) las casillas disponibles que tocan \(\mathcal J\), incluido \(c\), y \(E'\) las que sobreviven.
Para cada casilla \(i\), sea \(a_i\) el vector de incidencia de sus cuatro líneas.
Con \(h_f=1-\sum_{i\in f}p_i\) en las diagonales, pongamos
\begin{align*}
 S_0&=\sum_i p_i a_i a_i^{\mathsf T}+\operatorname{diag}(0_R,0_C,h_D,h_A),
 & R&=S_0^+,\\
 t_{ic}&=a_i^{\mathsf T}Ra_c,
 & w_i&=6p_i+f_i,\qquad f_i=\sum_{\{u,v\}\subset \familias(i)}R_{uv}.
\end{align*}
Las sumas sobre pares no ordenados usan dos líneas distintas. En \(E'\), escribamos
\(y_i=p_i^{(c)}-p_i\), donde \(p^{(c)}\) es el punto estacionario del hijo, y
\(q_i=y_i/p_i+t_{ic}\). La respuesta del hijo da \(q_i=O_b(m^{-2})\).
El potencial completo del prefactor es
\(V=\ell+\tfrac12\log\det S_{0,\rm gf}+3\sum_i p_i^2\). En los soportes conexos de
la clase, \(L=H+\ell+\tfrac12\log(2\pi)\) y \(\ell\) incluye el cofactor del predictor.
En \(\det S_{0,\rm gf}\) se usa el mismo calibre en padre e hijo.
La corrección lineal y su centrado se leen en
\[
 \delta_c^0=-w_c+\sum_i p_i t_{ic}w_i,\qquad
 \sum_{c\in r}p_ca_c=S_0 e_r,\qquad a_i^{\mathsf T}RS_0 e_r=(a_i)_r .
\]
La última igualdad usa que \(a_i\) es ortogonal al modo nulo fila--columna.
Por tanto \(\sum_{c\in r}p_c\delta_c^0=0\) para cada fila \(r\).

La consistencia por inserción~\cite[Teorema 7.1 y Corolario 7.2]{pc} da
\(L(B+c)-L(B)=\log[p_c(1+\delta_c)]+r_c^{\rm PC}\), con
\(\sum_{c\in r}p_c\delta_c=0\), \(\delta_c=O_b(m^{-1})\) y \(r_c^{\rm PC}=O_b(m^{-2})\).
Al desarrollar la exponencial en esos pesos de masa uno y después el logaritmo se obtiene
\(\log\rho(B,r)=\sum_{c\in r}p_c r_c^{\rm PC}+O_b(m^{-3})\).
El cambio de pesos en el término lineal cuesta \(O_b(m^{-3})\).

\paragraph{Qué términos deja el cambio de soporte.}
Descompongamos \(S_0\) en bloques \(B_0,C_0,C_0^{\mathsf T},D_0\) según \(I,\mathcal J\).
Sea \(U_0=\sum_{i\in F}p_i a_{i,I}a_{i,I}^{\mathsf T}\),
\(\mathcal M=B_0-U_0\), \(d_0=C_0\mathbf1\), \(\beta=Ra_c\),
\(z=\beta_I\), \(t=\beta_{\mathcal J}-\mathbf1\) y \(r=d_0+\mathcal M z=-C_0t-U_0z\).
El mínimo de la energía cuadrática del hijo es \(Q_*=d_0^{\mathsf T}\mathcal M^+d_0\).
Al insertar \(d_0=r-\mathcal Mz\) y usar \(S_0\beta=a_c\) se obtiene
\[
 Q_*=a_c^{\mathsf T}Ra_c-4+12p_c+E_Q,\qquad
 E_Q=z^{\mathsf T}U_0z-t^{\mathsf T}D_0t+r^{\mathsf T}\mathcal M^+r .
\]
En efecto \(D_0=(1-p_c)I+p_c\mathbf1\mathbf1^{\mathsf T}\) y
\(\mathbf1^{\mathsf T}D_0\mathbf1=4+12p_c\). Se tiene \(r=O_b(m^{-2})\) por entrada
y \(r^{\mathsf T}\mathcal M^+r=O_b(m^{-3})\) por la brecha y el número \(O_b(m)\) de líneas.

Sea \(\gamma_f=h_f^{(c)}-h_f\) el cambio de vacancia en una diagonal retenida.
Taylor de la entropía del apéndice~\ref{app:predictor}, con los pesos del hijo, da
\begin{align*}
 \Delta H-\log p_c&=\tfrac32\sum_{F\setminus\{c\}}p_i^2-\tfrac12 Q_*+A_3+O_b(m^{-3}),\\
 A_3&=\sum_{F\setminus\{c\}}p_i^3+\tfrac16\sum_{E'}p_i(y_i/p_i)^3
 +\tfrac32\sum_{E'}\frac{y_i^2}{1-p_i}+\tfrac16\sum_{f\in I_D}\gamma_f^3/h_f^2 .
\end{align*}
Aquí \(I_D\) recorre ambas familias diagonales retenidas. Los términos cuárticos cuestan
\(O_b(m^{-3})\) por \(\sum p_i=m\), \(p_i=O_b(m^{-1})\) y
\(|y_i/p_i|+|\gamma_f/h_f|=O_b(m^{-1})\). La ortogonalidad en el mínimo deja un error
cuadrático en las correcciones relativas \(O_b(m^{-2})\) a la respuesta lineal.

Para el cofactor escribamos \(\Pi_m=vv^{\mathsf T}/(2m)\), donde \(v\) vale \(1,-1,0,0\)
en las cuatro familias. Sean \(H_*=(R+\Pi_m)_{II}\), \(C_{\mathcal J}=D_0+(\Pi_m)_{\mathcal J\mathcal J}\),
\(V_{\mathcal J}=C_0+(\Pi_m)_{I\mathcal J}\) y \(\Pi'=v_Iv_I^{\mathsf T}/[2(m-1)]\).
En \(I\), \(S_*\) tiene entrada \(-p_i\) en cada par ordenado de líneas distintas
de \(i\in F\setminus\{c\}\). En forma análoga \(Y\) tiene entrada \(y_i\) en los pares de
\(i\in E'\). Sus restantes entradas son cero. El complemento de Schur da
\begin{align*}
 D_*&=S_*+Y+\Pi'-(\Pi_m)_{II}+V_{\mathcal J} C_{\mathcal J}^{-1}V_{\mathcal J}^{\mathsf T},\\
 Q_c&=\tfrac{m^2}4\{\operatorname{tr}[(H_*D_*)^2]-\operatorname{tr}(S_*^2)\},\\
 E_R&=\operatorname{tr}[(C_{\mathcal J}-I)(H_{\mathcal J}-I)],\qquad H_{\mathcal J}=(R+\Pi_m)_{\mathcal
 J\mathcal J}.
\end{align*}
Para la normalización del calibre sean \(S'_0\) la matriz del hijo y
\(\Delta_0=S'_0-(S_0)_{II}\). Su contribución se obtiene de
\[
 E_g=\log[m/(m-1)]-\log\det C_{\mathcal J}
   +\operatorname{tr}[(\Pi_m)_{II}\Delta_0]
   +\operatorname{tr}[H_*(\Pi'-(\Pi_m)_{II})].
\]
Como \(v_{\mathcal J}\perp\mathbf1\), \(\det C_{\mathcal J}=(1+3p_c)(1-p_c)^2(1-p_c+1/m)\).
Los logaritmos dan \(m^{-2}-p_c/m+6p_c^2+O_b(m^{-3})\). Las dos trazas valen
\(-(1-p_c)/m\) y \((m-1)/m^2+v_{\mathcal J}^{\mathsf T}R_{\mathcal J\mathcal J}v_{\mathcal J}/[2m(m-1)]\),
 respectivamente,
por \(S_0v=Rv=0\). Como \(R_{\mathcal J\mathcal J}=I+O_b(m^{-1})\), queda \(E_g=m^{-2}+6p_c^2+O_b(m^{-3})\).
El cubo del logaritmo del determinante deja \(C^{(3)}=-\operatorname{tr}(S_*^3)/6\).
En efecto, las filas de \(S_*\) tienen suma absoluta \(O_b(m^{-1})\);
las entradas de \(D_*-S_*\) son \(O_b(m^{-2})\) y las de \(H_*-I\) son \(O_b(m^{-1})\).
En la traza cúbica, un factor de \(H_*D_*-S_*\) cuesta así \(O_b(m^{-3})\) en total.
Las potencias de orden cuatro o superior tienen la misma cota por la norma de filas.
La traza pura \(\operatorname{tr}(S_*^2)/4\) aporta
\(3\sum_{F\setminus\{c\}}p_i^2/2\), pues hay seis pares ordenados en cada casilla retirada.

En la parte lineal de Schur, sea \(f_i^I\) la restricción de \(f_i\) a pares retenidos.
Las cuatro entradas diagonales de \(S_0R=I-\Pi_m\) sobre \(\mathcal J\) dan
\[
 \sum_{i\in F}p_i(f_i-f_i^I)
 =-\operatorname{tr}[(R-I)_{\mathcal J\mathcal J}]-1/m-p_cf_c .
\]
Al juntar la traza lineal, \(-Q_*/2\), las dos copias de
\(3\sum_{F\setminus\{c\}}p_i^2/2\) y \(-3\Delta\sum p_i^2\), los términos con \(w_i\)
son \(\sum_F p_iw_i-w_c-\sum_{E'}y_iw_i\) y queda \(p_cf_c-3p_c^2\).
Restar \(\delta_c^0\) convierte las dos sumas en las de la última línea de la fórmula
\begin{align*}
 L(B+c)-L(B)-\log p_c-\delta_c^0
   &=\Delta_c V+P_c+O_b(m^{-3}),\\
 P_c&=-E_g/2+Q_c/m^2+C^{(3)}+A_3-3\sum_{E'}y_i^2\\
 &\quad-3p_c^2+p_cf_c-(E_Q+E_R)/2\\
 &\quad+\sum_F p_i(1-t_{ic})w_i-\sum_{E'}p_iq_iw_i .
\end{align*}


Puesto que \(\delta_c-\delta_c^0=O_b(m^{-2})\), desarrollar \(\log(1+\delta_c)\) da
\[
 r_c^{\rm PC}=\Delta_c V+P_c+\tfrac12(\delta_c^0)^2
                   +(\delta_c^0-\delta_c)+O_b(m^{-3}).
\]
El último paréntesis desaparece exactamente al promediar en la fila. Los errores cúbicos
acumulados cuestan \(O_b(n\sum_{m>j}m^{-3})=O_b(n^{-1})\). La tabla identifica qué
parte de \(m^2[P_c+(\delta_c^0)^2/2]\) origina cada contribución al perfil.
\begin{center}\small
\begin{tabular}{ll}
\toprule
Término de inserción a escala \(m^2\) & Contribución al perfil\\ \midrule
\(-m^2 E_g/2\)& \(G_{\rm cof}\)\\
\(m^2 A_3\)&\(f2E+f2z+f2F\)\\
\(m^2 C^{(3)}\)&\(f2F+C_{3,\rm cuad}\)\\
\(-m^2(E_Q+E_R)/2\)&\(f3F+f3c\)\\
\(m^2(-3p_c^2+p_c f_c)\)&\(f4=f4_{\rm expl}+f4_{L0}\)\\
\(m^2\sum_F p_i(1-t_{ic})w_i\)&\(f5c+f5F\)\\
\(-m^2\sum_{E'}p_iq_iw_i\)&\(f6\)\\
\(m^2(\delta_c^0)^2/2\)&\(f7\)\\
primera parte de \(Q_c\)&\(Q_1+\Delta_1\)\\
segunda parte de \(Q_c\), con \(-3m^2\sum_{E'}y_i^2\)&\(Q_2+\Delta_2\)\\
\bottomrule
\end{tabular}
\end{center}
\subsection{Las dieciséis contracciones del perfil}\label{sec:contracciones}
Denotamos por \(\langle g\rangle_i=\int_{[0,1]^2}\nu_1(i)g(i)\,di\) la media de área y
por \(\langle\cdot\rangle_c\) la misma media en el centro. En la fibra \(\lambda\)
de \(c\), \(d\omega_\lambda(i)=\eta_\lambda\nu_1(i)\,dy\), con \(y\) la coordenada libre
de fila o columna y, en las diagonales, la coordenada de fila. Para una función \(g\),
\(\int_F g=\sum_{\lambda\in \familias(c)}\int_{i\in\lambda}g(i,\lambda)\,d\omega_\lambda(i)\).
La medida \(d\rho\) de una línea libre es \(dy\) en
cada familia axial y \(U_f\,dy/(2x)\) en cada hoja de una familia diagonal.
El factor \(1/2\) ya realiza la media en la hoja de esa línea libre. Los demás índices
enteros libres aportan sus propias clases de paridad, según su parametrización. En una fibra
axial el desplazamiento tiene dos paridades de igual peso. En una diagonal el desplazamiento
mantiene el color. Un cruce diagonal--antidiagonal sólo existe en hojas compatibles.

\paragraph{Qué disponibilidad lleva una intersección.}
En una configuración fijamos las líneas de sus casillas marcadas y las líneas que se integran
con \(d\rho\). Sea \(\mathcal A\) ese conjunto. A cada otra línea \(\ell\) de una configuración
finita se le asigna una variable de Bernoulli independiente \(X_\ell\) de media \(U_\ell\).
En \(\mathcal A\) se pone \(X_\ell=1\). Sean \(\familias(z)\) las cuatro líneas de \(z\),
\(Y_z=\prod_{\ell\in \familias(z)}X_\ell\) y \(\pesodirecto(z)=W(z)Y_z\). Un cruce ausente da
 \(\pesodirecto=0\).
Así el peso de varias casillas directas \(Z\) se obtiene de
\[
 \psi(Z)=\mathbb E\prod_{z\in Z}Y_z
 =\prod_{\ell\in(\bigcup_{z\in Z}\familias(z))\setminus\mathcal A}U_\ell .
\]
El mismo indicador se usa en todo el monomio. Para el núcleo \(L_0\) y la corrección
\(\theta\) de la sección~\ref{sec:profile-data}, la respuesta de dos líneas es
\[
 E(a,b)=
 \begin{cases}
 (L_0)_{aa}+\theta_a,&a=b,\\
 (L_0)_{ab}-\pesodirecto(a\cap b),&a\ne b .
 \end{cases}
\]
Una casilla marcada aporta \(\pesodirecto(i)=W(i)\). Abreviamos
\(z_f=\sum_{j\in \familias(c)}E(f,j)\), \(T_i=\sum_{a\in \familias(i)}z_a\),
\(t_j=z_j\) para \(j\in \familias(c)\), y \(w(i)=\sum_{\{a,b\}\subset \familias(i)}(L_0)_{ab}\).
En particular \(t_j\) y \(w(i)\) en una casilla marcada son deterministas.
Para \(i\) que sólo comparte \(\lambda\) con \(c\), sea
\(\sigma_i=\sum_{a\in \familias(i)\setminus\{\lambda\}}z_a\).

\paragraph{Cubo, energía y corrección de fila.}
Las contribuciones del cubo de la acción son
\begin{align*}
 f2E&=\left\langle-\tfrac16\mathbb E T_i^3+\tfrac32 W(i)\mathbb E T_i^2\right\rangle_{i,c},\\
 f2z&=-\tfrac16\left\langle\int_{D\cup A}(u_f/U_f)\,\mathbb E z_f^3\,d\rho(f)\right\rangle_c,\\
 f2F&=\left\langle\int_F W(i)^2\right\rangle_c .
\end{align*}
En \(f2E\) se marcan \(i,c\), independientes en área. En \(f2z\) se marcan \(c,f\).
Para verificar los momentos en estos dos casos, sumemos sobre los pares de líneas que definen
\(T_i\) o \(z_f\). Sean \(A_1,A_2,A_3\) las sumas respectivas de \(L_0-K\),
\(K(W_{ab}-K)\) y \(-K(W_{ab}-K)(W_{ab}-2K)\), con \(W_{ab}=W(a\cap b)\) o cero.
En posición general las transversales de cruces distintos son distintas. El segundo momento es
\(A_1^2+A_2\) y el tercero \(A_1^3+3A_1A_2+A_3\), por los momentos centrados de cada Bernoulli.
El factor \(u_f/U_f\) es la vacancia del punto estacionario en la línea retenida.

Con \(I_2=\langle W(c)^2\rangle_c\) y \(\Xi=\langle W(c)w(c)\rangle_c\), se obtiene
\[
 G_{\rm cof}=-\tfrac12-3I_2,\qquad f4=-9I_2+\Xi,\qquad f5c=-3\Xi .
\]
En efecto \(mf_c\) tiende al factor \(w(c)-6W(c)\), y \(1-t_{cc}=-3+O_b(m^{-1})\).
Las otras contracciones del centro y de la fibra son
\begin{align*}
 f3F&=-\tfrac12\left\langle\int_F\mathbb E\sigma_i^2\right\rangle_c,\\
 f3c&=\tfrac12\left\langle\sum_{j\in \familias(c)}t_j^2
 -\sum_{j,h\in \familias(c)}[W(c)(1-\delta_{jh})+\Pi_{jh}]
                         [E(h,j)+\Pi_{hj}]\right\rangle_c,\\
 f5F&=-\left\langle\int_F w(i)\mathbb E T_i\right\rangle_c,\\
 f7&=\tfrac12\left\langle\left[-w(c)+\int_F w(i)
                  +\left\langle w(i)\mathbb E T_i\right\rangle_i\right]^2\right\rangle_c .
\end{align*}
Aquí \(\Pi_{jh}=v_jv_h/2\) es el núcleo del perfil. El primer sumando de \(f3c\)
procede de \(t^{\mathsf T}D_0t\) y el segundo de \(E_R\).
En \(f3F\), \(i,c\) están marcadas en la misma fibra. Sean \(Z\) los cruces existentes
\(a\cap j\), con \(a\in \familias(i)\setminus\{\lambda\}\) y \(j\in \familias(c)\setminus\{\lambda\}\),
de familias distintas. En posición general son casillas distintas. Escribiendo
\(\psi_z=\psi(\{z\})\), el signo negativo de cada parte directa desaparece en la covarianza,
\begin{align*}
 \mathbb E\sigma_i^2={}&(\mathbb E\sigma_i)^2+\sum_{z\in Z}W(z)^2\psi_z(1-\psi_z)\\
 &+\sum_{\substack{z,z'\in Z\\z\ne z'}}W(z)W(z')
       [\psi(\{z,z'\})-\psi_z\psi_{z'}].
\end{align*}
Por ejemplo, en coordenadas relativas a \(c=(0,0)\), si \(i=(0,s)\), los cruces
\((s,s)\) y \((s,0)\) comparten la fila \(s\). Cuando existen dentro del tablero, su peso
conjunto es \(\psi_z\psi_{z'}/U_R\). En \(f7\) las dos copias de la suma tienen índices
libres distintos. Tras fijar \(c\), sus coincidencias accidentales tienen medida nula.
Así se cuadran las integrales de la respuesta centrada antes de promediar en \(c\).

\paragraph{Triángulos de líneas en la traza cúbica.}
Tres líneas del producto \(S_*^3\) concurren en una casilla borrada o cortan tres casillas
distintas que completan un cuadrado con \(c\). En el primer caso sus seis órdenes cancelan
el divisor \(6\), y el resultado es otra copia de \(f2F\).
Para el segundo, tomemos la base \(((1,0),(0,1))\) o \(((1,1),(1,-1))\) y multipliquemos
cada vector por un signo independiente, obteniendo \(v,w\). Sean
\(Z=\{c,c+hv,c+hw,c+h(v+w)\}\), \(h>0\), y \(n_{RC}\) el número de líneas axiales de \(Z\).
Las tres coordenadas libres son las de \(c\) y el paso primitivo \(h\). El resultado es
\[
 C_{3,\rm cuad}=\sum_{v,w}\int_{\{(c,h):\,Z\subset[0,1]^2\}}
  x^{n_{RC}-3}\prod_{z\in Z}W(z)
  \prod_{\ell\in(\bigcup_{z\in Z}\familias(z))\cap(D\cup A)}U_\ell\,dc\,dh .
\]
Los cuatro factores \(W\) incluyen el peso del centro. En la cuadratura el lado cero
lleva la mitad del valor límite del integrando. En la integral tiene medida cero.

\paragraph{Respuesta de segundo orden en el exterior.}
La fila \(f6\) requiere \(m^2q_i\), pues la respuesta de primer orden en \(y_i/p_i\)
se cancela con \(t_{ic}\). Para verificar la fuente de segundo orden, pongamos
\(X=(a_{i,I}^{\mathsf T})_{i\in E'}\), \(D_p=\operatorname{diag}(p_i)_{i\in E'}\).
Las ecuaciones estacionarias restadas dan \(y_i/p_i=(X\eta)_i+(B_1)_i\), donde
\(\eta_f=\gamma_f/h_f\) en las diagonales y las otras componentes son multiplicadores
de fila y columna. Con \(\theta_i=y_i/p_i\), la diferencia exacta de los logaritmos
de las ecuaciones de estacionariedad da
\[
 (B_1)_i=\theta_i-\log(1+\theta_i)-3\log(1-y_i/(1-p_i))
 +\sum_{f\in \familias(i)\cap(D\cup A)}[\log(1+\eta_f)-\eta_f].
\]
Las ecuaciones de masa dan
\(\mathcal M\eta=d_0-X^{\mathsf T}D_pB_1\). Por tanto
\[
 q=B_1+X\mathcal M^+(r-X^{\mathsf T}D_p B_1).
\]
Taylor en las casillas y las vacancias da \(m^2B_1=L+O_b(m^{-1})\), con
\[
 L_i=\tfrac12(mt_{ic})^2-\frac{3mp_i(mt_{ic})}{1-p_i}
       -\tfrac12\sum_{f\in \familias(i)\cap(D\cup A)}(mz_f^{\rm fis})^2,\qquad
 z^{\rm fis}=(Ra_c)_I .
\]
Aquí se han usado \(y_i/p_i=-t_{ic}+O_b(m^{-2})\) y
\(\gamma_f/h_f=-z_f^{\rm fis}+O_b(m^{-2})\).
En el complemento del modo nulo, la identidad \(S_0R=I-\Pi_m\) por bloques da
\(\|\mathcal M^+-R_{II}\|_\infty=O_b(m^{-1})\) en norma de suma por fila,
pues las líneas retenidas pierden masa \(O_b(m^{-1})\) y la inversa está acotada.
Así \(m^2q=L+XR_{II}(m^2r-X^{\mathsf T}D_pL)+O_b(m^{-1})\).

Expresamos la contracción en los mismos núcleos del perfil. Fijemos el centro \(c\) y sea
\(\Lambda(i)=T_i^2/2-3W(i)T_i-\tfrac12\sum_{f\in \familias(i)\cap(D\cup A)}z_f^2\).
Para una línea retenida \(u\), escribamos
\begin{align*}
 \omega(u)={}&-\sum_{j\in \familias(c)} \pesodirecto(u\cap j)
     \left[t_j+\sum_{a\in \familias(u\cap j)\setminus\{j\}}z_a\right]
      -\int_{i'\in u}\mathbb E\Lambda(i')\,d\omega_u(i'),\\
 q(i)={}&\Lambda(i)+\sum_{u\in \familias(i)}
 \left[\omega(u)+\int E(u,u')\omega(u')\,d\rho(u')\right],\\
 f6={}&-\left\langle w(i)\mathbb E q(i)\right\rangle_{i,c}.
\end{align*}
Se omiten los cruces inexistentes en la primera suma. El término de fibra integrado en
\(\omega\) es la media con \(i'\) marcada y \(c\) fijo. En los productos de la suma sobre
cruces se usa \(\psi\) para todo el monomio. Las dos partes de \(\omega\) provienen de
\(m^2r=-m^2(C_0t+U_0z)\) y de \(-X^{\mathsf T}D_pL\). La identidad y la parte integral
de \(R_{II}\) producen los dos términos entre corchetes de \(q(i)\). La media de área
lleva el signo de \(-m^2\sum_{E'}p_iq_iw_i\). El Lema~\ref{lem:f6} demuestra su
reordenación por trazas, incluido el recorrido de la figura~\ref{fig:rectangulos}.

\begin{lema}[Reordenación de la contribución \(f6\)]\label{lem:f6}
En la malla de lado \(s\), sean \(\mathcal I{}\) la incidencia casilla--línea de la
sección~\ref{sec:rectangulos}, \(\beta_e=\nu_1(e)/s^2\) y
\(N_v=\mathcal I{}^{\mathsf T}\operatorname{diag}(v)\mathcal I{}\). Los productos de vectores de casillas
son puntuales. Con \(K,L_0,\rho,w\) de la sección~\ref{sec:profile-data}, definamos
\begin{align*}
 D&=L_0-K,& V_{ij}&=K_{ij}(W_{ij}-K_{ij}),&
 \Psi&=\rho^{-1}\mathcal I{}^{\mathsf T}(\beta w),\\
 Z&=\Psi+D\rho\Psi,& h&=w-\mathcal I{}Z,& \alpha&=\beta h,\\
 a&=\mathcal I{}^{\mathsf T}\alpha,& b_\beta&=\mathcal I{}^{\mathsf T}\beta,&
 M&=\mathcal I{}D\mathcal I{}^{\mathsf T}.
\end{align*}
Aquí \(W_{ij}\) es el peso del cruce o cero si éste no existe.
Escribamos \(U(e,j)=\sum_{u\in \familias(e)} K_{uj}-3W(e)\) y
\(t(c,e;j)=M_{ce}+\theta_j+U(c,j)+U(e,j)\).
La cuadratura del perfil físico, con el valor genérico en coincidencias accidentales,
da exactamente \(f6_s=-E_1+E_2\), donde
\begin{align*}
 E_1={}&\tfrac12\operatorname{tr}(N_\beta D N_\alpha D)
       +\tfrac12 b_\beta^{\mathsf T}Va-3b_\beta^{\mathsf T}D\mathcal I{}^{\mathsf T}(\alpha W)\\
     &-\tfrac12\sum_{f\in D_{\rm obl}}a_f
         [(D N_\beta D)_{ff}+(Vb_\beta)_f],\\
 E_2={}&\sum_c\beta_c\sum_{j\in \familias(c)}\sum_{e\in j}
     \frac{\eta_j\nu_1(e)}s\,t(c,e;j)
     \sum_{u\in \familias(e)\setminus\{j\}}Z_u .
\end{align*}
En esta fórmula \(D_{\rm obl}\) reúne las dos familias oblicuas.
\end{lema}
\begin{proof}
Fijado \(c\), la respuesta de segundo orden consta de la fuente local \(\Lambda\) y
la fuente de líneas \(\omega\). Al sumar la respuesta contra \(\beta_e w(e)\), cada
\(u\in \familias(e)\) recibe \([\mathcal I{}^{\mathsf T}(\beta w)]_u=(\rho\Psi)_u\).
La parte integral de la respuesta añade \(\rho D\rho\Psi\) por simetría. La masa
total es, pues, \(\rho Z\). La fuente negativa de \(\omega\) integra \(\Lambda\) en
la fibra de \(u\). En la malla \(\rho_u\eta_u=1/s\), por lo que esa suma sustrae
\(\beta_e(\mathcal I{}Z)_e\) del peso de \(\Lambda(e;c)\). Queda la masa \(\alpha_e\).

Para dos marcas genéricas \(c,e\), la media de \(T_e\) es \(M_{ce}\).
Cada cruce directo tiene varianza \(V_{uj}\), pues su indicador al cuadrado es el propio
indicador. Los cruces distintos tienen transversales distintas en posición genérica.
Así \(\mathbb E T_e^2=M_{ce}^2+\sum_{u\in \familias(e),j\in \familias(c)}V_{uj}\).
Para \(f\in \familias(e)\) la misma cuenta da
\(\mathbb E z_f^2=(\sum_{j\in \familias(c)}D_{fj})^2+\sum_{j\in \familias(c)}V_{fj}\).
Al sumar \(\beta_c\alpha_e\mathbb E\Lambda(e;c)\), el cuadrado \(M_{ce}^2\)
da la traza de \(N_\beta D N_\alpha D\) por multiplicación de índices.
La varianza da \(b_\beta^{\mathsf T}Va\), y el momento lineal con \(W(e)\)
da \(b_\beta^{\mathsf T}D\mathcal I{}^{\mathsf T}(\alpha W)\).
Los cuadrados de \(z_f\) dan \((DN_\beta D)_{ff}+(Vb_\beta)_f\), uno por línea
oblicua de \(e\). Los coeficientes de \(\Lambda\) producen \(E_1\).
En \(N_v\) se incluyen las entradas diagonales \(\sum_{e\in u}v_e\), también cuando
el recorrido vuelve por la misma línea.

En la suma de cruces de \(\omega\), fijemos \(e=u\cap j\), con \(j\in \familias(c)\).
El peso directo por la medida de \(u\) es \(\rho_uK_{uj}=\eta_j\nu_1(e)/s\).
Condicionar el cruce a estar disponible marca entonces \(c,e\) en la fibra \(j\).
En la respuesta media \(M_{ce}\) se repone la diagonal \(\theta_j\).
Los tres cruces con \(j\) en cada extremo llevan ahora el peso cierto \(W\) de la
casilla marcada. Cambiar cada antiguo término \(-K\) a \(-W\) añade \(U(c,j)+U(e,j)\).
Queda \(t(c,e;j)\), multiplicado por \(Z_u\) para cada transversal \(u\ne j\) de \(e\).
La fuente de cruces tiene signo negativo en \(\omega\), y la definición de \(f6\)
tiene otro signo negativo. De ahí \(+E_2\) y \(-E_1\). Esta igualdad de sumas
finitas deja el error espacial de la extrapolación en el estado \status{N}.
\end{proof}

\paragraph{Las dos trazas del término mixto.}
Pongamos \(D=H_*D_*-S_*\). La ciclicidad de la traza da la identidad en el padre
\[
 Q_c=\tfrac12\sum_{u,v}(S_*)_{uv}(m^2D)_{vu}
       +\tfrac1{4m^2}\sum_{u,v}(m^2D)_{uv}(m^2D)_{vu}.
\]
Cada casilla de \(F\setminus\{c\}\) tiene tres líneas retenidas, que dan seis pares
ordenados en la primera suma. En la segunda se integran dos líneas libres.
Para expresar \(m^2D\) en el perfil escribamos \(\mathcal K(v,w)=E(v,w)+\Pi_{vw}\),
\(\mathcal V(v,j)=\pesodirecto(v\cap j)+\Pi_{vj}\), \(j\in \familias(c)\). Para líneas retenidas \(v,u\),
sea \(\mathcal Y(v,u)=-\pesodirecto(z)T_z\) si se cortan en \(z\in E'\), y cero en otro caso.
En esta fórmula \(T_z\) usa los mismos indicadores que \(\pesodirecto(z)\).
Finalmente, para \(w\ne u\) retenidas, sea \(\mathcal S(w,u)=-\pesodirecto(i')\) cuando
\(w\cap u=i'\in F\setminus\{c\}\), y cero en los demás pares. Definamos
\begin{align*}
 \mathcal B(v,u)&=\mathcal Y(v,u)+\Pi_{vu}
                    +\sum_{j\in \familias(c)}\mathcal V(v,j)\mathcal V(u,j),\\
 \mathcal E_2(v,u)&=\mathcal B(v,u)+\sum_w\mathcal K(v,w)\mathcal S(w,u)
                     +\int\mathcal K(v,w)\mathcal B(w,u)\,d\rho(w).
\end{align*}
Estas tres piezas proceden, con \(B_*=D_*-S_*\) y \(K_*=H_*-I\), de
\(m^2D=m^2B_*+(mK_*)(mS_*)+m^{-1}(mK_*)(m^2B_*)\).
En la suma discreta se toma cada \(i'=u\cap j\), \(j\in \familias(c)\), y cada una de
sus dos líneas transversales \(w\). Su disponibilidad está en \(\pesodirecto(i')\).
En la integral sí se marca \(w\), pues su disponibilidad ya está en \(d\rho(w)\).

Escribamos \(R(i)=\{(u,v):u,v\in \familias(i)\setminus\{\lambda\},\ u\ne v\}\) para una
casilla de la fibra \(\lambda\). Aplicando \(\psi\) a los monomios se obtienen
\begin{align*}
 L_1&=-\tfrac12\left\langle\int_F
          \sum_{(u,v)\in R(i)}\mathbb E\mathcal E_2(v,u)\right\rangle_c,\\
 L_2^{\rm int}&=\tfrac14\left\langle\iint
 \mathbb E\big[\mathcal E_2(u,v)\mathcal E_2(v,u)-\mathcal Y(u,v)\mathcal Y(v,u)\big]
                 \,d\rho(u)\,d\rho(v)\right\rangle_c .
\end{align*}
Las dos copias de una línea integrada internamente tienen variables libres separadas.
El bloque \(\mathcal Y\mathcal Y\) se sustrae antes de integrar. En el tablero cada
\(i\in E'\) aporta doce pares ordenados al bloque \(YY\), de modo que su contribución
es \(+3m^2\sum_{E'}y_i^2\), exactamente la opuesta a la de \(P_c\).
También se puede comprobar en el continuo. Para cada par de familias, la medida de los
cruces válidos por su disponibilidad es \(U_DU_A\,di\), con las hojas y el jacobiano del
cruce incluidos. Multiplicando por \(W(i)^2\) resulta \(\nu_1(i)W(i)\).

\paragraph{Cómo extraer \(\Delta_1\) y \(\Delta_2\).}
Las filas tabuladas son \(Q_1=L_1-\Delta_1\) y \(Q_2=L_2^{\rm int}-\Delta_2\).
Damos las correcciones por separado para poder reconstruir cada entrada. En todas las
sumas siguientes un cruce debe existir en el cuadrado y en la clase de paridad indicada.
Una igualdad de líneas se exige como identidad en los parámetros libres de esa clase.
Las coincidencias accidentales forman conjuntos de medida nula.

Para \(\Delta_1\), fijemos \(c,i,\lambda,(u,v)\in R(i)\) como en \(L_1\), y una
línea libre \(w\) integrada con \(d\rho(w)\). Pongamos \(z=v\cap w,\ p=w\cap u\).
Para cada \(a\in \familias(p)\setminus\{w,u\}\) y \(j\in \familias(c)\), pongamos \(t=a\cap j\).
Se marcan \(\familias(c)\cup \familias(i)\cup\{w\}\).
El monomio directo de \(\mathcal K(v,w)\mathcal Y(w,u)\) aporta la diferencia
\[
 \Delta_1=\tfrac12\left\langle\int_F\sum_{(u,v)\in R(i)}
 \int\sum_{a,j} W(z)W(p)W(t)
 \big[\psi(\{z,p,t\})-\psi(\{z\})\psi(\{p,t\})\big]\,d\rho(w)\right\rangle_c .
\]
Se toman los tres cruces distintos, con \(p\in E'\). La diferencia cuenta la transversal
común de \(z,t\) sólo una vez. Si es \(\ell\), el corchete vale
\(\psi(\{z\})\psi(\{p,t\})(U_\ell^{-1}-1)\).
Las tres partes directas llevan signo negativo en el producto de la primera traza
antes de aplicar su prefactor \(-1/2\).
Por ejemplo, en coordenadas relativas, para \(c=(0,0)\), \(i=(0,s)\), \(u=C_s\),
\(v=D_{-s}\), \(w=R_h\), \(a=A_{h+s}\), \(j=R_0\), resultan \(z=(h,h+s)\),
\(p=(h,s)\), \(t=(0,h+s)\). La columna \(h+s\) aparece en \(z\) y \(t\).

Para \(\Delta_2\), marquemos \(\familias(c)\cup\{u,v\}\), con \(u,v\) libres. Cuando existe el cruce,
el peso de sus transversales es
\[
 K_{ab}=W(a\cap b)\prod_{\ell\in \familias(a\cap b)\setminus\{a,b\}}U_\ell .
\]
Para cada \(j\in \familias(c)\) tomamos \(p_2=v\cap j\),
\(w_2\in \familias(p_2)\setminus\{v,j\}\), \(t_2=u\cap w_2\), \(k_2=K_{vj}K_{uw_2}\).
Los dos patrones posibles para los otros dos vértices dan
\[
\begin{array}{c|c|c|c}
 &p_1,\ t_1 &k_1&\text{multiplicidad }s\\ \hline
 \mathcal K\mathcal S\ \text{por}\ \mathcal K\mathcal S&
 u\cap j,\ v\cap w_1 & K_{uj}K_{vw_1}&1\\
 \mathcal Y\ \text{por}\ \mathcal K\mathcal S&
 v\cap u,\ a\cap j & K_{vu}K_{aj}&2
\end{array}
\]
En la primera fila \(w_1\) recorre \(\familias(p_1)\setminus\{u,j\}\). En la segunda \(a\)
recorre \(\familias(p_1)\setminus\{v,u\}\). En cada patrón se conservan cuatro vértices distintos.
Para cada pareja de cruces que define \(k_d\), \(d=1,2\), reunamos las transversales de
sus dos cruces, es decir, las líneas cuyos factores \(U\) aparecen en \(k_d\).
Si los dos conjuntos comparten una línea \(\ell\) fuera de las marcadas, se incluye el
sumando siguiente. Los cuatro cruces forman entonces un cuadrado y la transversal es única.
\[
 \Delta_2=\tfrac14\left\langle\iint\sum_{\text{patrones admitidos}}
          s\,k_1k_2(U_\ell^{-1}-1)\,d\rho(u)\,d\rho(v)\right\rangle_c .
\]
Se puede obtener la lista desarrollando los productos de \(\mathcal E_2\).
Un peso \(\pesodirecto\) aporta las cuatro coordenadas de su casilla menos las líneas ya marcadas.
Para comparar dos etiquetas de una misma familia se restan sus expresiones afines en
las coordenadas libres. Sólo una diferencia idénticamente nula aporta el factor
\(U_\ell^{-1}-1\). Para hacer la comprobación basta comparar el término constante y
los coeficientes de cada parámetro en las cuatro etiquetas \(r,q,r-q,r+q\). En
\(\Delta_1\) se obtienen 24 elecciones de familias y 80 al distinguir las clases admisibles.
En \(\Delta_2\) quedan 24 elecciones en cada fila de la tabla. Por ejemplo, en la primera,
con \(j=R_0,u=A_s,v=D_t,w_1=C_s,w_2=C_{-t}\), los dos vértices exteriores
\((s+t,s),(s+t,-t)\) comparten la fila \(s+t\). En la parte mixta, intercambiar las copias da el mismo
integrando con \(u,v\) permutadas, de donde procede el \(2\) de la tabla.
Finalmente,
\begin{align*}
 d(x)={}&G_{\rm cof}+f2E+f2z+2f2F+f3F+f3c+f4+f5c+f5F\\
        &+f6+f7+C_{3,\rm cuad}+Q_1+\Delta_1+Q_2+\Delta_2 .
\end{align*}


\subsection{Límite uniforme de las medias y suma sobre los niveles}
La regularidad antes de la parada da \(p_i\le C_b/m\), una vacancia diagonal
minorada y una brecha de filas y columnas~\cite[\S7]{pc}.
En consecuencia \(c_b I\le S_0+\Pi_m\le 4I\) para algún \(c_b>0\).
Para ver la cota inferior, la forma cuadrática es la suma de
\(\sum_i p_i(a_i^{\mathsf T}y)^2\) y \(\sum_f h_f y_f^2\), más la proyección de calibre.
El segundo término controla las componentes diagonales y la brecha controla las axiales
ortogonales al calibre. La cota superior sigue de la diagonal unidad de \(S_0\), de que su
masa fuera de la diagonal por fila es a lo sumo tres y de que la proyección actúa en su núcleo.

Sea \(A_m=S_0+\Pi_m-I\), \(M=(I+A_m)^{-1}=R+\Pi_m\) y \(r_b=1-c_b/4<1\).
El polinomio \(P_N=\tfrac14\sum_{h=0}^N[I-(I+A_m)/4]^h\) satisface
\(\|M-P_N\|_{\rm op}\le r_b^{N+1}/c_b\). Para controlar una entrada multiplicada
por \(m\) usamos \(M=I-A_m+A_mMA_m\). La suma absoluta de una fila de \(A_m\) es
a lo sumo cuatro y su máximo es \(O_b(m^{-1})\), luego
\(\|A_me_u\|_2=O_b(m^{-1/2})\) para todo \(u\). Así
\[
 \sup_{u,v}m\,|[A_m(M-P_N)A_m]_{uv}|\le C_b r_b^{N+1}.
\]
Esto acota el error en cada entrada de \(m(R-I)\) al sustituir \(M\) por \(P_N\)
en la identidad anterior. Las entradas originales y sus truncaciones están acotadas.

A \(N\) fijo, las potencias de \(S_0-I\) son sumas sobre sucesiones de líneas adyacentes.
Las inserciones de \(\Pi_m\) se simplifican usando \((S_0-I)\Pi_m=-\Pi_m\).
Un producto de recorridos es una suma finita de tipos de incidencia. Si dos etiquetas afines
son idénticas en un tipo, se cuenta la línea una vez. Una coincidencia adicional impone
una ecuación no constante y deja \(O_{b,N}(n^{d-1})\) posiciones entre las \(O(n^d)\)
de un tipo con \(d\) parámetros. En la diagonal, la ida y vuelta de longitud dos fuerza
la misma casilla en ambos pasos. Su contribución es \(3m\sum_{i\in u}p_i^2\), pues hay tres
elecciones de la otra línea. Su exceso sobre el producto de las medias da precisamente
\(\theta_u=3\eta_u s_u-\int K_{uv}^2\,d\rho(v)\). Los cruces entre diagonales se
cuentan en las clases de paridad de la sección~\ref{sec:contracciones}.

A cada tipo fijo se aplica el seguimiento de líneas hasta el corte auxiliar \(b/4\)
de~\cite[\S4]{tc}. A tolerancia fija \(\eta>0\), fuera de una excepción de probabilidad
exponencial, \(mp_i=W(i)(1+O_b(\eta+m^{-1}))\) en los sitios vivos y
\(h_f=u_f/U_f+O_b(\eta+m^{-1})\). Para un número fijo de marcas el riesgo en un paso es
la suma de sus riesgos individuales salvo \(O_{b,N}(n^{-2})\) por dobles ocupaciones.
Al multiplicar estas supervivencias hasta el nivel \(m\), una marca axial aporta \(x\)
y una diagonal aporta \(U_f(x)\), pues \(\partial_x\log U_f=(1-u_f)/U_f\).
La aproximación es uniforme en las marcas y en \(j\le m\le n\). Se toma primero \(n\to\infty\)
y luego \(\eta\downarrow0\). Las sumas normalizadas de cada clase son sumas de Riemann
sobre politopos fijos. Los factores del perfil son continuos en \(x\) con módulo común en
\([b,1]\) en cada pieza. Por tanto los truncados convergen uniformemente en los niveles
a integrales continuas de los recorridos.

Para aplicar el mismo argumento a \(mV\), sea \(\kappa_i=a_i^{\mathsf T}Ra_i-4\).
La expansión del prefactor~\cite[\S7]{pc} da, con la última suma sobre todas las líneas,
\[
 V=-14\sum_i p_i^3-\tfrac32\sum_i p_i^2\kappa_i
       +\tfrac94\sum_u\left(\sum_{i\in u}p_i^2\right)^2+O_b(m^{-2}).
\]
Las doce entradas fuera de la diagonal en \(m\kappa_i\) aportan \(-12W(i)\) en la
casilla marcada. Combinadas con \(-14\) dan \(4\langle W^2\rangle_i\).
Las partes largas y la diagonal corregida dan los sumandos con \(L_0,\theta\) de
la sección~\ref{sec:profile-data}. El cuadrado de la suma de fibra da su último término.

En las contracciones, el centro tiene masa uno, el exterior masa a lo sumo uno y las
cuatro fibras masa a lo sumo cuatro. La medida de las líneas libres tiene masa \(\le6/b\).
Los pasos de cuadrados tienen masa total \(\le8/b\). Al desarrollar la diferencia de productos,
la cota de cada entrada se multiplica sólo por masas y factores acotados para \(b\) fijo.
Así la cola del grado \(N\) es \(\le C_b r_b^{N+1}\), uniforme en el nivel tanto para la media
física como para su integral límite. El resto de la inserción en la escala \(m^2\) es
\(O_b(m^{-1})\). Tomar primero \(n\to\infty\) y después \(N\to\infty\) da
\begin{align*}
 \sup_{j<m\le n}\Big|
 \mathbb E_\nu\!\left[m^2\sum_{c\in i_m}p_c
       \big(P_c+(\delta_c^0)^2/2\big);H_j\right]-d(m/n)\Big|&\longrightarrow0,\\
 \sup_{j\le m\le n}\left|\mathbb E_\nu[mV(B_m);H_j]-\mathcal V(m/n)\right|
  &\longrightarrow0 .
\end{align*}
Aquí \(i_m\) es la fila revelada y \(H_j\) la llegada regular corta. Para promediar la fila
se usa primero el indicador de supervivencia anterior a su elección. Dado el padre, la fila
es uniforme y la media se convierte en \(m^{-1}\sum_c p_c(\cdot)\).
Cambiar ese indicador por \(\mathbf1_{H_j}\) en los observables acotados cuesta
\(O_b(\nu(H_j^c))\). La segunda fórmula usa todas las familias en el potencial.
Este argumento identifica una media sobre padres y da continuidad del límite en \([b,1]\).
La evaluación decimal de esas integrales por mallas requiere además una cota de su error espacial.

\paragraph{Los extremos y el cambio de ley.}
El error uniforme \(\epsilon_n(b)\) de la primera media contribuye a lo sumo
\(\epsilon_n(b)n\sum_{m>j}m^{-2}=o_b(1)\). El término de \(P_c+(\delta_c^0)^2/2\)
da por tanto \(\int_b^1d(x)x^{-2}\,dx\). Para \(\Delta_cV\), la probabilidad de un hijo
dada la fila bajo \(\nu\) es \(p_c(1+O_b(m^{-1}))\), y \(|\Delta_cV|\le C_bm^{-2}\).
La sustitución en la media cuesta \(O_b(m^{-3})\). Los incrementos centrados hasta la parada
forman una martingala de segundo momento \(O_b(\sum_{m>j}m^{-4})\).
Si \(M_\tau\) es su suma detenida, las diferencias son ortogonales en \(L^2(\nu)\).
Así \(n|\mathbb E_\nu[M_\tau;H_j]|\le n(\mathbb E_\nu M_\tau^2)^{1/2}\nu(H_j)^{1/2}\).
Como \(\sum_{m>j}m^{-4}=O(j^{-3})\), el coste a escala \(n\) es \(O_b(n^{-1/2})\).
En cada historia de \(H_j\) los incrementos realizados suman \(V(B_j)-V(B_n)\).
La segunda media uniforme da \(\mathcal V(b)/b-\mathcal V(1)\).

Hasta la parada, \(n|A|\le C_b\) para el predictor ordinario. Por el Lema~\ref{lem:previa},
\(n|R_Q(n)|\) también está acotado. En \(H_j\) ponemos
\(s=A-R_Q(n)+g_T+\Gamma\). Allí \(|s|=O_b(n^{-1})\) y la densidad inversa de historias
es \(e^{-s-g^\flat}\). Por \eqref{eq:arrival-estimate},
\[
 \int_{H_j}|d\mu-d\nu|
 \le\mathbb E_\mu[|1-e^{-s}|+e^{-s}|1-e^{-g^\flat}|;H_j]=O_b(n^{-1}) .
\]
Con las estimaciones de salida del apéndice~\ref{app:signed} se sigue que
\(n\mathbb E_\mu A=n\mathbb E_\nu[A;H_j]+o_b(1)\).
Para las torres pongamos \(u_e=g_T(M_e^{\rm fis})-g_T(M)\).
La máscara \(M_e^{\rm fis}\) se obtiene del menor simple \(M_e\) aplicando los dos vetos
diagonales del hijo. El cambio es \(O_b(m^{-2})\), incluido el paso que entra en la salida.
Así \(\log\rho^\sharp-\log\rho=\log\mathbb E_{\nu,i}e^{u_e}
 =\mathbb E_{\nu,i}u_e+O_b(m^{-4})\).
La densidad detenida y las cotas ponderadas de salida dan
\(D(\mu^\tau\Vert\nu^\tau)=\mathbb E_\mu[A+\log X]-R_Q(n)=O_b(1)\).
Sean \(d_m\) las entropías condicionales de las dos leyes de hijo dada la fila,
iguales a cero después de la parada. Las filas tienen la misma ley uniforme.
La descomposición de la entropía relativa en sus incrementos condicionales da
\(\sum_m\mathbb E_\mu d_m=D(\mu^\tau\Vert\nu^\tau)\). La desigualdad de
Pinsker~\cite{ct} acota la diferencia de las dos medias de \(u_e\) por \(2C_bm^{-2}\sqrt{d_m/2}\).
Por Jensen y Cauchy--Schwarz, la suma de estas diferencias a escala \(n\) cuesta
\(O_b(n(\sum_{m>j}m^{-4})^{1/2}D(\mu^\tau\Vert\nu^\tau)^{1/2})=O_b(n^{-1/2})\).
Los incrementos realizados de torres se telescopan y añaden
\((F(b)-1/12)/b-(F(1)-1/12)\) a los extremos anteriores. En conclusión, a cada corte fijo,
\[
 \lim_{n\to\infty} n\mathbb E_\mu A^\sharp
 =\int_b^1\frac{d(x)}{x^2}\,dx+\frac{G(b)}b-G(1)
 =\int_b^1\frac{\numerr(x)}{x^2}\,dx .
\]

